
    %%%%%%%%%%%%%%%%%%%%%%%%%%%%%%%%%%%%%%%%%%%%%%%%%%%%%%%%%%%%%%%%%%%%%%%%%%%%%%%%
    %%% Classe do documento
    \documentclass[12pt, addpoints]{exam}
    \UseRawInputEncoding

    %%%%%%%%%%%%%%%%%%%%%%%%%%%%%%%%%%%%%%%%%%%%%%%%%%%%%%%%%%%%%%%%%%%%%%%%%%%%%%%%
    % preambulo.tex
    %
    % Arquivo de configuração de packages para uso o LaTeX, conforme minhas
    % preferências e modelos pessoais.
    %
    % Para maiores informações, visite:
    %    
    %
    % NÃO ALTERE SE NÃO SOUBER O QUE ESTÁ FAZENDO!
    %%%%%%%%%%%%%%%%%%%%%%%%%%%%%%%%%%%%%%%%%%%%%%%%%%%%%%%%%%%%%%%%%%%%%%%%%%%%%%%%


    %%%%%%%%%%%%%%%%%%%%%%%%%%%%%%%%%%%%%%%%%%%%%%%%%%%%%%%%%%%%%%%%%%%%%%%%%%%%%%%%
    %%% Estruturas de controle:
    \input{static/utils/controles}

    %%%%%%%%%%%%%%%%%%%%%%%%%%%%%%%%%%%%%%%%%%%%%%%%%%%%%%%%%%%%%%%%%%%%%%%%%%%%%%%%
    %%% Compilação condicional em PDF:
    \input{static/utils/pdf}

    %%%%%%%%%%%%%%%%%%%%%%%%%%%%%%%%%%%%%%%%%%%%%%%%%%%%%%%%%%%%%%%%%%%%%%%%%%%%%%%%
    %%% Configurações de layout da página:
    \input{static/utils/layout}

    %%%%%%%%%%%%%%%%%%%%%%%%%%%%%%%%%%%%%%%%%%%%%%%%%%%%%%%%%%%%%%%%%%%%%%%%%%%%%%%%
    %%% Configurações para autores:
    \input{static/utils/autores}

    %%%%%%%%%%%%%%%%%%%%%%%%%%%%%%%%%%%%%%%%%%%%%%%%%%%%%%%%%%%%%%%%%%%%%%%%%%%%%%%%
    %%% Cabeçalho e rodapé:
    %%% Atenção! É MUITO DIFÍCIL ajustar apropriadamente os running headers
    %%% e running footers da primeira vez. Você pode confiar no LaTeX e deixar que
    %%% ele faça o ajuste automaticamente ou pode quebrar a cabeça e
    %%% tentar melhorar algo que o LaTeX já faz muito bem, mesmo que não fique
    %%% exatamente do jeito que você quer. O que você prefere? Se quiser ajustar
    %%% manualmente, descomente a linha abaixo e faça as configurações no arquivo
    %%% "utils/headings.tex".
    %\input{static/utils/headings}

    %%%%%%%%%%%%%%%%%%%%%%%%%%%%%%%%%%%%%%%%%%%%%%%%%%%%%%%%%%%%%%%%%%%%%%%%%%%%%%%%
    %%% Configuração para as fontes do documento:
    %%% Se você quiser usar fontes próprias ou do sistema, precisará ajustar as
    %%% configurações no arquivo "utils/fontes.tex".
    %%% ATENÇÃO: por padrão eu utilizo XeLaTeX com fontes proprietárias projetadas
    %%% por Matthew Butterick, adquiridas em https://mbtype.com, que não podem ser
    %%% distribuídas devido à licença de utilização. Se você usar XeLaTeX, você
    %%% DEVERÁ OBRIGATORIAMENTE ajustar as fontes no arquivo "utils/fontes.tex".
    \input{static/utils/fontes}

    %%%%%%%%%%%%%%%%%%%%%%%%%%%%%%%%%%%%%%%%%%%%%%%%%%%%%%%%%%%%%%%%%%%%%%%%%%%%%%%%
    %%% Sumário:
    \input{static/utils/toc}

    %%%%%%%%%%%%%%%%%%%%%%%%%%%%%%%%%%%%%%%%%%%%%%%%%%%%%%%%%%%%%%%%%%%%%%%%%%%%%%%%
    %%% Matemática:
    \input{static/utils/matematica}

    %%%%%%%%%%%%%%%%%%%%%%%%%%%%%%%%%%%%%%%%%%%%%%%%%%%%%%%%%%%%%%%%%%%%%%%%%%%%%%%%
    %%% Notas de rodapé:
    \input{static/utils/footnotes}

    %%%%%%%%%%%%%%%%%%%%%%%%%%%%%%%%%%%%%%%%%%%%%%%%%%%%%%%%%%%%%%%%%%%%%%%%%%%%%%%%
    %%% Referências cruzadas, links, citações:
    \input{static/utils/crossrefs}

    %%%%%%%%%%%%%%%%%%%%%%%%%%%%%%%%%%%%%%%%%%%%%%%%%%%%%%%%%%%%%%%%%%%%%%%%%%%%%%%%
    %%% Configurações para os hiperlinks em PDF:
    \input{static/utils/pdfconfig}

    %%%%%%%%%%%%%%%%%%%%%%%%%%%%%%%%%%%%%%%%%%%%%%%%%%%%%%%%%%%%%%%%%%%%%%%%%%%%%%%%
    %%% Glossário e índice remissivo:
    \input{static/utils/glossindex}

    %%%%%%%%%%%%%%%%%%%%%%%%%%%%%%%%%%%%%%%%%%%%%%%%%%%%%%%%%%%%%%%%%%%%%%%%%%%%%%%%
    %%% Referências bibliográficas:
    \input{static/utils/bibliografia}

    %%%%%%%%%%%%%%%%%%%%%%%%%%%%%%%%%%%%%%%%%%%%%%%%%%%%%%%%%%%%%%%%%%%%%%%%%%%%%%%%
    %%% Cores:
    \input{static/utils/cores}

    %%%%%%%%%%%%%%%%%%%%%%%%%%%%%%%%%%%%%%%%%%%%%%%%%%%%%%%%%%%%%%%%%%%%%%%%%%%%%%%%
    %%% Computação:
    \input{static/utils/computacao}

    %%%%%%%%%%%%%%%%%%%%%%%%%%%%%%%%%%%%%%%%%%%%%%%%%%%%%%%%%%%%%%%%%%%%%%%%%%%%%%%%
    %%% Gráficos:
    \input{static/utils/graficos}

    %%%%%%%%%%%%%%%%%%%%%%%%%%%%%%%%%%%%%%%%%%%%%%%%%%%%%%%%%%%%%%%%%%%%%%%%%%%%%%%%
    %%% Ícones:
    \input{static/utils/icones}

    %%%%%%%%%%%%%%%%%%%%%%%%%%%%%%%%%%%%%%%%%%%%%%%%%%%%%%%%%%%%%%%%%%%%%%%%%%%%%%%%
    %%% Blocos:
    \input{static/utils/blocos}

    %%%%%%%%%%%%%%%%%%%%%%%%%%%%%%%%%%%%%%%%%%%%%%%%%%%%%%%%%%%%%%%%%%%%%%%%%%%%%%%%
    %%% Boxes coloridos:
    \input{static/utils/colorbox}

    %%%%%%%%%%%%%%%%%%%%%%%%%%%%%%%%%%%%%%%%%%%%%%%%%%%%%%%%%%%%%%%%%%%%%%%%%%%%%%%%
    %%% Tabelas:
    \input{static/utils/tabelas}

    %%%%%%%%%%%%%%%%%%%%%%%%%%%%%%%%%%%%%%%%%%%%%%%%%%%%%%%%%%%%%%%%%%%%%%%%%%%%%%%%
    %%% Ambientes de listas:
    \input{static/utils/listas}

    %%%%%%%%%%%%%%%%%%%%%%%%%%%%%%%%%%%%%%%%%%%%%%%%%%%%%%%%%%%%%%%%%%%%%%%%%%%%%%%%
    %%% Ambientes floats e similares:
    \input{static/utils/floats}

    %%%%%%%%%%%%%%%%%%%%%%%%%%%%%%%%%%%%%%%%%%%%%%%%%%%%%%%%%%%%%%%%%%%%%%%%%%%%%%%%
    %%% Ajustes para a classe EXAM:
    \input{static/utils/exam}

    %%%%%%%%%%%%%%%%%%%%%%%%%%%%%%%%%%%%%%%%%%%%%%%%%%%%%%%%%%%%%%%%%%%%%%%%%%%%%%%%
    %%% Ajustes para textos:
    \input{static/utils/texto}

    %%%%%%%%%%%%%%%%%%%%%%%%%%%%%%%%%%%%%%%%%%%%%%%%%%%%%%%%%%%%%%%%%%%%%%%%%%%%%%%%
    %%% Ambientes, aliases, comandos e customizações em geral para serem
    %%% utilizadas nos documentos. Defina aqui qualquer customização específica
    %%% para seus documentos!
    \input{static/utils/customizacoes}

    %%%%%%%%%%%%%%%%%%%%%%%%%%%%%%%%%%%%%%%%%%%%%%%%%%%%%%%%%%%%%%%%%%%%%%%%%%%%%%%%
    %%% Ajuste do layout, espaçamento de linhas e etc.:
    \geometry{a4paper, portrait, left=2.0cm, right=2.0cm, top=2.0cm, bottom=2.0cm}
    \singlespacing

    %%%%%%%%%%%%%%%%%%%%%%%%%%%%%%%%%%%%%%%%%%%%%%%%%%%%%%%%%%%%%%%%%%%%%%%%%%%%%%%%
    %%% Configurações de cabeçalho e rodapé (não use fancyhdr pois ocorre conflito):
    \pagestyle{headandfoot}
    \runningheadrule
    \firstpageheader{}{}{}
    \runningheader{Sem disciplina}{}{Avaliação}
    \firstpagefooter{}{}{Página \thepage\ de \numpages}
    \runningfooter{}{}{Página \thepage\ de \numpages}
    \runningfootrule

    %%%%%%%%%%%%%%%%%%%%%%%%%%%%%%%%%%%%%%%%%%%%%%%%%%%%%%%%%%%%%%%%%%%%%%%%%%%%%%%%
    %%% Configurações para as propriedades do PDF:
    \hypersetup{
    hidelinks,           % Comente para web, descomente para impressão
    colorlinks=true,      % True para web, False para impressão
    pdftitle=Prova Pre Enade: Avaliação,
    pdfauthor=Professor,
    pdfsubject=Sem disciplina,
    pdfkeywords=['Sem tag'],
    pdfinfo={
        CreationDate={}, % Ex.: D:AAAAMMDDHH24MISS
        ModDate={}       % Ex.: D:AAAAMMDDHH24MISS
    }
    }
    \input{static/utils/pdfconfig}

    %%%%%%%%%%%%%%%%%%%%%%%%%%%%%%%%%%%%%%%%%%%%%%%%%%%%%%%%%%%%%%%%%%%%%%%%%%%%%%%%
    %%% Começa o documento
    \begin{document}

    %%%%%%%%%%%%%%%%%%%%%%%%%%%%%%%%%%%%%%%%%%%%%%%%%%%%%%%%%%%%%%%%%%%%%%%%%%%%%%%%
    %%% Folha de instruções padronizadas da para a prova
    \pdfbookmark[1]{Instruções}{instrucoes}
    \begin{figure}[H]
        \begin{center}
            \includegraphics[scale=0.3]{{static/imgs/logo-.png}}
        \end{center}	
    \end{figure}

    \vspace{-1cm}

    \begin{table}[H]
        \centering
        \begin{tabular}{|llll|}
            \hline
            \multicolumn{2}{|l|}{\textbf{Disciplina:} Sem disciplina} & 
            \multicolumn{1}{l|}{\multirow{3}{*}{\textbf{\makecell{Nota:\\ \,\\ \,}}}} & 
            \multirow{3}{*}{\textbf{\makecell{Coordenador:\\ \,\\ \,}}} \\ 
            \cline{1-2}
            \multicolumn{2}{|l|}{\textbf{Professor:} Professor \hspace{4.72cm} } & 
            \multicolumn{1}{l|}{} & \\ 
            \cline{1-2}
            \multicolumn{2}{|l|}{\textbf{Aluno:}} & 
            \multicolumn{1}{l|}{} & \\ 
            \hline
            \multicolumn{1}{|l|}{\textbf{Turma:} \hspace{6.3cm} } & 
            \multicolumn{1}{l|}{\textbf{Semestre:}} & 
            \multicolumn{2}{l|}{\textbf{Valor:} 10 pontos} \\ 
            \hline
            \multicolumn{1}{|l|}{\textbf{Data:}} & 
            \multicolumn{3}{l|}{\textbf{Avaliação:}} \\ 
            \hline
        \end{tabular}
    \end{table}

    \begin{center}
        \fbox{\setlength\fboxsep{0.3cm} \fbox{\parbox{15.4cm}{
            \begin{center}
                \textbf{INSTRUÇÕES PARA A REALIZAÇÃO DESTA AVALIAÇÃO:}
            \end{center}
            \begin{itemize}[leftmargin=*]
                \item Esta é a primeira avaliação da disciplina Sem disciplina, 
                    e engloba o conteúdo referente Sem tag.
                \item Esta avaliação contém 20 questões objetivas, \textbf{todas obrigatórias}.
                \item \textbf{Leia atentamente cada questão!} As questões objetivas terão uma,
                    e apenas uma única, resposta a ser marcada.
                \item \textbf{Desligue o celular} antes de começar. A avaliação é
                    \textbf{individual} e \textbf{sem consulta}.
                \item As questões podem ser respondidas, na prova, com lápis ou caneta. Ao final
                    da prova \textbf{{VOCÊ DEVERÁ PREENCHER O GABARITO COM CANETA
                    DE COR PRETA}} \textbf{\underline{OU AZUL}} para que seja feita a correção
                    automatizada através da leitura do padrão de respostas marcado no
                    gabarito.
                \item \textbf{CUIDADO ao preencher o gabarito} pois eles são \textbf{INDIVIDUAIS
                    e NOMEADOS} (cada estudante tem um gabarito próprio específico, com um
                    código de identificação único). \textbf{Questões rasuradas são anuladas}
                    pelo software de leitura do gabarito.
                \item Ao final da prova, \textbf{DEVOLVA A PROVA E O GABARITO} para o professor.
                \item Siga todas as normas de \textbf{Integridade Acadêmica} da disciplina.
                    Alunos que forem flagrados com qualquer espécie de ``cola'' ou trocando
                    informações com outros alunos terão suas avaliações recolhidas, as notas
                    zeradas, e a situação será encaminhada para a coordenação para a aplicação
                    das penalidades previstas pela universidade.
                \item Com 90 minutos de avaliação você terá tempo de até 4,min por questão,
                    em média.
                \item Boa avaliação!
            \end{itemize}
            \vspace{2cm}
        }}}
    \end{center}
    \newpage

    %%%%%%%%%%%%%%%%%%%%%%%%%%%%%%%%%%%%%%%%%%%%%%%%%%%%%%%%%%%%%%%%%%%%%%%%%%%%%%%%
    %%% Inicia o bloco de questões:
    \begin{questions}

    \newpage
    \question

Um compilador detecta:
\begin{choices}
\choice erros de sintaxe do programa 
\choice erros que podem ocorrer durante a execução do programa
\choice erros nos resultados gerados pelo programa
\choice erros aritméticos
\choice todos os erros citados acima
\end{choices}

\question

Quais das seguintes afirmações são verdadeiras? As Métricas de software servem para:



\begin{itemize}

    \item[I)] indicar a qualidade do produto e avaliar a produtividade.

    \item[II)] auxiliar na melhoria do processo.

    \item[III)] formar uma base para as estimativas e justificar a aquisição de ferramentas.

    \item[IV)] determinar se a utilização de um método traz benefícios ou não.

\end{itemize}
\begin{choices}
\choice Apenas as alternativas I, II e IV.
\choice Todas as alternativas.
\choice Nenhuma delas.
\choice Apenas as alternativas II e III.
\choice Apenas as alternativas I, IV.
\end{choices}

\question

Qual o valor do atributo \( E.val \) após a análise da expressão "4 / 2 / 2" para o esquema de tradução a seguir?



\begin{itemize}

    \item \( E \rightarrow T / E1 \ \{ E.val = T.val / E1.val \} \)

    \item \( E \rightarrow T \ \{ E.val = T.val \} \)

    \item \( T \rightarrow \text{digito} \ \{ T.val = \text{val(digito)} \} \)

\end{itemize}


\begin{choices}
\choice 3
\choice 4
\choice 8
\choice 1
\choice 2
\end{choices}

\question

Qual das seguintes condições não é necessária para a ocorrência de um \textit{deadlock}?


\begin{choices}
\choice Alocação parcial de recursos a processos.
\choice Escalonamento preemptivo de recursos.
\choice Haver compartilhamento de recursos por processos.
\choice Processos em espera circular.
\choice Uso mutuamente exclusivo de recursos por processos.
\end{choices}

\question

O número de \textit{strings binárias} de comprimento 7 e contendo um par de zeros consecutivos é


\begin{choices}
\choice 91
\choice 95
\choice 90
\choice 92
\choice 94
\end{choices}

\question

Considere o seguinte código para implementar exclusão mútua entre dois processos \(i\) e \(j\): 

Processo \(P_i\)

\begin{lstlisting}[language=C, basicstyle=\ttfamily\small]

do {

    while (turn != i); 

    // secao critica

    secao critica

    turn = j;

    // saída da secao critica

    codigo restante

} while (1);

\end{lstlisting} 
\begin{choices}
\choice A implementação garante exclusão mútua.
\choice Os processos fazem espera ativa.
\choice Exige alternância estrita. 
\choice A implementação garante progresso. 
\choice Um processo bloqueia o outro mesmo não estando na seção crítica.
\end{choices}

\question

Considere as seguintes afirmativas sobre as linguagens usadas para análise sintática:

\begin{enumerate}[label=\Roman*.]

    \item A classe LL(1) não aceita linguagens com produções que apresentem recursões diretas à esquerda (ex. \(L \rightarrow La\)) mas aceita linguagens com recursões indiretas (ex. \(L \rightarrow Ra, R \rightarrow Lb\)).

    \item A linguagem LR(1) reconhece a mesma classe de linguagens que LALR(1).

    \item A linguagem SLR(1) reconhece uma classe de linguagens maior que LR(0).

\end{enumerate}

Selecione a afirmativa correta:
\begin{choices}
\choice As afirmativas I e II são verdadeiras
\choice Apenas a afirmativa III é verdadeira
\choice As afirmativas I e III são verdadeiras
\choice As afirmativas II e III são verdadeiras
\choice As afirmativas I e III são falsas
\end{choices}

\question

Considere as seguintes afirmações:

\begin{enumerate}[label=\Roman*.]

    \item A classe de linguagens aceita por um Autômato Finito Determinístico (AFD) não é a mesma que um Autômato Finito Não Determinístico (AFND).

    \item Para algumas expressões regulares não é possível construir um AFD.

    \item A expressão regular \((b + ba)^+\) aceita as "strings" de b's e a's começando com b e não tendo dois a's consecutivos.

\end{enumerate}

Selecione a afirmativa correta:
\begin{choices}
\choice As afirmativas I e II são verdadeiras.
\choice As afirmativas I e III são falsas.
\choice As afirmativas II e III são falsas.
\choice As afirmativas I e III são verdadeiras.
\choice Apenas a afirmativa III é verdadeira.
\end{choices}

\question

(ENADE 2017) No modelo relacional de dados, uma junção entre tabelas cria uma

outra tabela derivada de duas ou mais tabelas de acordo com os critérios

especificados para a junção, de modo similar às regras da teoria dos

conjuntos. Nos conjuntos a seguir, as áreas hachuradas representam os resultados

das consultas SQL em que ``id' é uma chave primária e ``fk' é uma chave

estrangeira:

\begin{figure}[H]

  \centering

  \includegraphics[width=0.5\linewidth]{static/imgs/defaul.png}

\end{figure}

Assinale a opção em que a declaração SQL é equivalente ao conjunto apresentado a

seguir:

\begin{figure}[H]

  \centering

  \includegraphics[width=0.5\linewidth]{static/imgs/defaul.png}

\end{figure}


\begin{choices}
\choice \begin{verbatim}SELECT *

FROM tabela1 t1

RIGHT JOIN tabela2 t2 ON (t1.id = t2.fk)

WHERE t2.fk <> NULL;

\end{verbatim}


\choice \begin{verbatim}SELECT *

FROM tabela1 t1

INNER JOIN tabela2 t2 ON (t1.id = t2.fk);

\end{verbatim}


\choice \begin{verbatim}SELECT *

FROM tabela1 t1

WHERE NOT EXISTS (SELECT 1

                  FROM tabela1 t2

                  WHERE t2.id = tk.fk);

\end{verbatim}


\choice \begin{verbatim}SELECT *

FROM tabela1 t1 WHERE EXISTS (SELECT 1

                              FROM tabela2 t2

                              WHERE t2.id = t1.fk);

\end{verbatim}


\choice \begin{verbatim}SELECT *

FROM tabela1 t1

RIGHT JOIN tabela2 t2 ON (t1.id = t2.fk)

WHERE t1.id IS NULL;

\end{verbatim}


\end{choices}

\question

Qual das seguintes afirmações sobre crescimento assintótico de funções não é verdadeira:
\begin{choices}
\choice Se \( f(n) = O(g(n)) \) então \( g(n) = O(f(n)) \)
\choice Se \( f(n) = O(g(n)) \) e \( g(n) = O(h(n)) \) então \( f(n) = O(h(n)) \)
\choice \( \log n^2 = O(\log n) \)
\choice \( 2n^2 + 3n + 1 = O(n^2) \)
\choice \( 2^{n+1} = O(2^n) \)
\end{choices}

\question

Se \( \lim_{n \to \infty} \left( \frac{1}{n^2} + \frac{2}{n^2} + \cdots + \frac{n-1}{n^2} \right) = L \) então


\begin{choices}
\choice \( L = 2 \)
\choice \( L = 0 \)
\choice \( L = \infty \)
\choice \( L = 1/2 \)
\choice \( L = 1 \)
\end{choices}

\question

Considere as seguintes afirmações sobre resolução de problemas em IA.

\begin{enumerate}[label=\Roman*.]

    \item A* é um conhecido algoritmo de busca heurística.

    \item O Minimax é um dos principais algoritmos para jogos de dois jogadores, como o xadrez.

    \item Busca em espaço de estados é uma das formas de resolução de problemas em IA.

\end{enumerate}

São corretas:
\begin{choices}
\choice I, II e III
\choice Apenas III
\choice Apenas I e II
\choice Apenas II e III
\choice Apenas I e III
\end{choices}

\question

Pode-se afirmar que o gráfico da função \(y = 2 + \frac{1}{x - 1}\) é o gráfico da função \(y = \frac{1}{x}\):
\begin{choices}
\choice nenhuma das anteriores.
\choice transladado uma unidade para a direita e duas unidades para baixo
\choice transladado uma unidade para a esquerda e duas unidades para cima
\choice transladado uma unidade para a esquerda e duas unidades para baixo
\choice transladado uma unidade para a direita e duas unidades para cima
\end{choices}

\question

Qual dos protocolos abaixo pode ser caracterizado como protocolo de roteamento do tipo estado de enlace?
\begin{choices}
\choice ICMP
\choice BGP-4
\choice IGMP
\choice RIP2
\choice OSPF
\end{choices}

\question

Um banco de dados (BD) é uma coleção logicamente coerente de dados relacionados

e persistentes. Em relação às características de um banco de dados, assinale a

alternativa verdadeira:
\begin{choices}
\choice Representa um aspecto (ou parte) do mundo real, parte essa chamada de ``catálogo' ou ``dicionário de dados``.
\choice Não precisa de dados atuais e verdadeiros para ser útil (para representar corretamente a realidade atual).
\choice Armazena dois grandes ``tipos`` de dados: os dados do usuário final e os metadados, sendo que os metadados ficam armazenados nas mesmas tabelas que os dados dos usuários e isso nos permite saber a estrutura dos bancos de dados.
\choice Não precisam ser obrigatoriamente computadorizados, existem bancos de dados manuais e/ou em papel.
\choice Geralmente é projetado e usado para finalidades gerais, sendo extremamente genéricos e adaptáveis para qualquer situação.
\end{choices}

\question

Para que serve a segmentação de um processador (\textit{pipelining})?
\begin{choices}
\choice permitir a execução de mais de uma instrução por ciclo de relógio
\choice simplificar o conjunto de instruções
\choice aumentar a velocidade do relógio
\choice reduzir o número de instruções estáticas nos programas
\choice simplificar a implementação do processador
\end{choices}

\question

Professor Mac Sperto propôs o seguinte algoritmo de ordenação, chamado de Super Merge, similar ao \textit{merge sort}: divida o vetor em 4 partes do mesmo tamanho (ao invés de 2, como é feito no \textit{merge sort}). Ordene recursivamente cada uma das partes e depois intercale-as por um procedimento semelhante ao procedimento de intercalação do \textit{merge sort}. Qual das alternativas abaixo é verdadeira?
\begin{choices}
\choice Super Merge está correto, mas consome tempo \( O(\text{merge sort}) \)
\choice Super Merge não está correto. Não é possível ordenar quebrando o vetor em 4 partes
\choice Super Merge está correto, mas consome tempo menor que \( O(\text{merge sort}) \)
\choice Super Merge está correto, mas consome tempo maior que \( O(\text{merge sort}) \)
\choice Nenhuma das afirmações acima está correta
\end{choices}

\question

Suponha que \( T \) seja uma árvore AVL inicialmente vazia, e considere a inserção dos elementos \(10, 20, 30, 5, 15, 2\) em \( T \), nesta ordem. Qual das sequências abaixo corresponde a um percurso de \( T \) em pré-ordem:
\begin{choices}
\choice 20, 10, 5, 2, 15, 30
\choice 10, 5, 2, 20, 15, 30
\choice 2, 5, 10, 15, 20, 30
\choice 30, 20, 15, 10, 5, 2
\choice 15, 10, 5, 2, 20, 30
\end{choices}

\question

Considere \( C(x) \) uma função que define a complexidade de um problema \( x \); \( E(x) \) uma função que define o esforço (em termos de tempo) exigido para se resolver o problema \( x \). Sejam dois problemas denominados \( p1 \) e \( p2 \). Assinale a alternativa correta.


\begin{choices}
\choice \( E(p1 + p2) < E(p1) + E(p2) \)
\choice Nenhuma das alternativas anteriores
\choice Se \( C(p1) < C(p2) \) então \( E(p1) > E(p2) \)
\choice Se \( C(p1) < C(p2) \) então \( E(p1) < E(p2) \)
\choice \( C(p1 + p2) < C(p1) + C(p2) \)
\end{choices}

\question

(Adaptado do ENADE 2005) Avalie o seguite diagrama relacional de um banco de dados:

\begin{figure}[H]

  \centering

  \caption{Banco de dados de peças e fornecedores}

  \label{fig:pecas}

  %\fbox{

    \includegraphics[width=0.5\linewidth]{static/imgs/defaul.png}

  %}\\

  %\footnotesize{Fonte: xxxx.}

\end{figure}

Assinale a opção que apresenta um comando SQL que permite obter uma lista que

contenha o nome de cada fornecedor que tenha fornecido alguma peça, o código da

peça fornecida, a descrição dessa peça e a quantidade fornecida da referida

peça.
\begin{choices}
\choice \begin{verbatim}SELECT * FROM pecas INNER JOIN fornecimentos f ON (p.codigo = f.cod_peca) INNER JOIN fornecedores f2 ON (f.cod_forn = f2.cod_forn) ORDER BY f2.nome;\end{verbatim} 
\choice \begin{verbatim}SELECT f2.nome, p.codigo, p.descricao, f.quantidade FROM pecas INNER JOIN fornecimentos f ON (p.codigo = f.cod_peca) INNER JOIN fornecedores f2 ON (f.cod_forn = f2.cod_forn);\end{verbatim} 
\choice \begin{verbatim}SELECT * FROM pecas INNER JOIN fornecedores INNER JOIN fornecimentos;\end{verbatim} 
\choice \begin{verbatim}SELECT nome, codigo, descricao, quantidade FROM pecas INNER JOIN FORNECEDORES INNER JOIN FORNECIMENTOS;\end{verbatim} 
\choice \begin{verbatim}SELECT DISTINCT p.nome, p.codigo, p.descricao, f.quantidade FROM pecas p INNER JOIN fornecimentos f ON (p.codigo = f.cod_peca);\end{verbatim} 
\end{choices}



    %%%%%%%%%%%%%%%%%%%%%%%%%%%%%%%%%%%%%%%%%%%%%%%%%%%%%%%%%%%%%%%%%%%%%%%%%%%%%%%%
    %%% Finaliza o bloco de questões:
    \end{questions}
    \end{document}
    